\documentclass[11pt]{article}

\usepackage[a4paper,margin=1in]{geometry}
\usepackage[T1]{fontenc}
\usepackage[utf8]{inputenc}
\usepackage{lmodern}
\usepackage{amsmath,amssymb}
\usepackage{physics}
\usepackage{hyperref}
\usepackage{enumitem}
\usepackage{setspace}
\usepackage{mdframed}

\hypersetup{colorlinks=true,linkcolor=blue,citecolor=blue,urlcolor=blue}
\setstretch{1.1}

\newmdenv[linewidth=1.1pt,innertopmargin=5pt,innerbottommargin=5pt,%
  innerleftmargin=8pt,innerrightmargin=8pt]{resultbox}

\title{Lessons from a closed programme\\[2pt]
\large A post-mortem of Saturated Superfluid Vacuum theory}
\author{Stig Norland}
\date{2026-06-16}

\begin{document}
\maketitle

\begin{abstract}
The Saturated Superfluid Vacuum (SSV) programme --- an attempt to derive the
particle, force, gravity, and cosmology sectors from a single logarithmic
Schr\"odinger order parameter with the vacuum treated as a real superfluid medium
--- is closed. No surviving route delivers the quantities it set out to derive.
This is a post-mortem, written to the programme's own standing rule that only
correctly negative results are proof. Its purpose is to extract what generalises.
The recurring diagnosis is \emph{form yes, magnitude no}: the framework reproduces
the kinematic \emph{forms} of analogue and emergent gravity --- weak-field time
dilation, area-law horizon entropy, the Unruh/Hawking temperature, a logarithmic
flat-rotation tail --- essentially for free, because those follow from locality
and symmetry, while every dynamical \emph{magnitude} (Newton's constant $G$, the
cosmological scale, the particle masses) had to be conceded as an input. The most
durable outputs are two negatives: a structural no-go --- a scalar/acoustic vacuum
cannot carry a spin-2 graviton or tensor gravitational-wave polarisations --- and
a concrete mode-counting demonstration that a literal holographic screen
undershoots $G$ by roughly 38 orders of magnitude. We also record the methodology
that made clean failure possible, and the honest home of the surviving ontology
(foundations / philosophy of physics). The lessons transfer beyond SSV to
emergent-gravity and superfluid-vacuum model-building generally.
\end{abstract}

\section{What this document is}

A research programme can end in two ways: it can be abandoned, or it can be
\emph{closed} --- worked until the reasons it cannot succeed are themselves
results. SSV is closed in the second sense. The point of writing the ending down
is that the reasons are more transferable than the framework was. Almost
everything of value SSV produced is either a re-derivation of results already
owned by the analogue- and emergent-gravity literature, or a negative; and the
negatives are sharp enough to be useful to anyone building a vacuum out of a
condensate. This note collects them, states the failures at full strength, and
draws the lessons. It supersedes nothing; it concludes.

\section{The programme, in brief}

SSV proposed that the physical vacuum is a saturated superfluid described by a
macroscopic order parameter $\Psi=\sqrt{\rho}\,e^{i\theta}$, that the fundamental
particles are topological defects (vortices, breathers) in $\Psi$, and that the
four forces, gravity, and cosmology all follow from one master dynamics --- a
logarithmic Schr\"odinger equation (LogSE) with a chiral-shear term and a
saturation potential. The ambition was unusually total: from $m_e$ and $\alpha$
alone, to derive the mass spectrum, the force sector, Newton's constant, the
arrow of time, black-hole structure, galactic rotation curves, and the
cosmological constant. The logarithmic nonlinearity is itself not new --- it is
the Bialynicki-Birula--Mycielski equation, developed as a superfluid vacuum by
Zloshchastiev --- and the superfluid-vacuum lineage runs through Volovik. SSV's
distinctive bet was the totality: one equation, two inputs, everything else
derived.

\section{The method that let it fail cleanly}

Before the failures, the part worth keeping. The programme ran under a discipline
that is, in retrospect, its most exportable product:

\begin{itemize}[leftmargin=1.4em]
\item \textbf{A negative-results rule.} Only correctly negative results count as
  proof; positive results are merely suggestive and are never softened or buried.
  This single rule is what kept the programme from drifting into the permanent
  ``suggestive'' state that absorbs most heterodox physics.
\item \textbf{Pre-registration.} Each hypothesis and its decision rule were fixed
  in a tracked issue \emph{before} the computation ran, so a result could not be
  reinterpreted into a success after the fact.
\item \textbf{Claim-status tracking.} Every quantitative claim carried an explicit
  status --- \emph{derived}, \emph{coincidence}, or \emph{falsified} --- and a
  change of status had to propagate to the prose, the tables, and the record in
  one pass.
\item \textbf{Computation-backed analysis.} Analytic claims were backed by a
  tested script plus a result note, so a claimed derivation had an executable
  witness and a reference could not silently break.
\end{itemize}

\noindent This machinery is why the programme could kill its own central claims
rather than defend them. The lesson generalises: the value of pre-registration
and a hard negative-results rule is highest precisely where external peer
pressure is lowest.

\section{What was reproduced: the form}

In its gravity sector SSV reproduced, correctly, a sequence of standard results.
Identifying local time with the order-parameter phase rate via the Josephson
relation $\hbar\dot\theta=-\mu$ gives an emergent lapse, and weak-field this
yields $g_{00}=1+2\Phi/c^2$ with the Einstein-1911 (time-delay) half of light
deflection. The acoustic horizon --- the surface where the background flow turns
sonic --- behaves as a black-hole horizon with a definite surface gravity, no
interior singularity (the finite vortex-core size cuts off collapse), and an
Unruh/Hawking temperature of the right form. Quantising the LogSE Bogoliubov
fluctuations gives an area-law horizon entanglement entropy. A logarithmic
circulation-halo potential $\Phi_h=v_h^2\ln r$ produces a flat rotation-curve
tail, and ordinary Newtonian self-gravity of the baryons grows spiral arms and a
bar by swing amplification, with no central black hole and no dark-matter
particle.

Every item in that list is a known result of analogue/emergent gravity
(Unruh, Visser, Volovik, Jacobson) or of standard galactic dynamics. SSV
re-derived them inside one medium, which has pedagogical value, but none is new.
This matters for the diagnosis: the \emph{forms} came easily.

\section{What failed: magnitudes and structural no-gos}

\subsection{The particle masses became inputs and coincidences}

The mass sector reduced to two inputs, $m_e$ and $\alpha$, rather than deriving a
spectrum. The ``$\alpha$-harmonic mass ladder'' was demoted to a numerical
coincidence, not a derived spectrum. The muon as a derived breather mode failed
three independent null tests and was closed. The $W$ mass \emph{scale}
$m_W\sim m_e/\alpha^2$ is structurally derived, but the prefactor that brings it
to the measured value is an $\mathcal{O}(1)$ coincidence. The honest tally: the
sector reproduces scales it is handed and fits the rest.

\subsection{The scalar-gravity no-go (the durable negative)}

The sharpest structural result is a no-go. SSV's gapless carrier is a single
scalar phonon; the emergent acoustic metric is conformal-in-space plus a vector
flow term, with \emph{no} transverse-traceless part. So the framework's
gravitational radiation is helicity-0, the tensor overlap with a quadrupole
source is zero, and the isotropic part cancels in a differential readout. The
attempt to rescue a tensor mode from the spinor sector fails for a representation-
theoretic reason: $\bar{\mathbf 2}\otimes\mathbf 2=\mathbf 1\oplus\mathbf 3$
contains no spin-2. The same square-root-of-a-vector fact that gives the photon
spin-1 \emph{forbids} a spin-2 graviton. SSV's gravity is therefore
\emph{structurally scalar}: it can mimic light bending in the conformal/trace
sector but cannot be general relativity, and its predicted gravitational-wave
polarisation content is wrong (the original ``effective spin-2'' claim was
falsified and withdrawn; the surviving model passes the Hulse--Taylor test only
because a scalar emits no dipole). This is the Visser kinematics/dynamics wall,
stated explicitly: an acoustic metric supplies geometry, not Einstein dynamics.
It is a clean, transferable constraint on \emph{any} scalar/superfluid emergent-
gravity model.

\subsection{Newton's constant: a screen too sparse by 38 orders}

Entropic and holographic gravity (Jacobson, Verlinde, Padmanabhan) keep the
holographic screen abstract: one \emph{assumes} an area law with the Planck-scale
coefficient. SSV's distinctive move was to make the screen literal --- a real
acoustic horizon in a real medium --- so that the coefficient becomes computable
rather than assumed. The form survived: the horizon entanglement entropy is
area-law, computed by the correlator method (Peschel; Casini--Huerta) via the
Srednicki radial decomposition, with the chiral term exactly silent. But this is
the cheap half: the area law follows from the LogSE quadratic Hamiltonian being
\emph{local}, generic for any local free field (Srednicki; Bombelli--Koul--Lee--
Sorkin; Jacobson's entanglement-equilibrium). It \emph{earns} an area law the
papers had merely asserted; it does not surprise.

The coefficient is the test, and it fails. Closing the Jacobson construction with
the computed $dS=\eta\,dA$ gives $\nabla^2\Phi=4\pi G_{\mathrm{eff}}\,e/c^2$ with
$G_{\mathrm{eff}}=1/(4\hbar c\,\eta)$. Reproducing the measured $G$ requires
\begin{equation}
  N=\frac{1}{\alpha_G}=\frac{\hbar c}{G m_p^2}
   =\left(\frac{m_P}{m_p}\right)^2\simeq 1.7\times 10^{38}
\end{equation}
microstates per grain --- the Planck-to-proton hierarchy. The substrate supplies
$\mathcal{O}(10)$ localised Bogoliubov modes per healing-length cell, a few bits.
The literal screen is short by about $10^{37}$.

\begin{resultbox}
\noindent When a holographic screen is given a real microscopic substrate and its
localised degrees of freedom are \emph{counted}, the area law emerges for free but
the Planck-scale coefficient does not: the count falls short by $\sim$38 orders.
$G$ must enter below the effective theory's cutoff --- a trans-Planckian input,
not an emergent output.
\end{resultbox}

\noindent This is the Sakharov induced-gravity problem and the Susskind--Uglum
species problem, which Volovik names as the structural weak point of superfluid
gravity --- here \emph{demonstrated by direct mode-counting} rather than argued.
(Guard: the identity $(a_p/\ell_P)^2=1/\alpha_G$ merely restates the hierarchy and
must not be read backwards as deriving $N$; the negative is the \emph{measured}
count of $\mathcal{O}(10)$ confronted with the target.)

\subsection{The arrow of time: a coarse-graining, not a resolution}

SSV originally read its irreversibility as something new. A pre-registered
reversal-echo computation through a topological-change event settled it the other
way: the master dynamics is conservative and exactly time-reversal symmetric, and
exact reversal recovers the pre-event state (fidelity deficit $1-F\approx
8.7\times10^{-12}$ in 2D annihilation, $F=0.999997$ for 3D Hopf-link
reconnection, the deficit shrinking under refinement). The account of
irreversibility is therefore a \emph{coarse-graining} one in Boltzmann's lineage,
with the forgotten variables named concretely (phonon phases, Kelvin waves,
vortex tangle, correlations); it is not a new resolution of Loschmidt's paradox,
and the programme withdrew any such claim. The genuinely open half --- the
low-entropy past --- was never supplied. (The reconnection-reversibility
computation does survive as a stand-alone result for the quantum-turbulence
literature, shorn of the arrow-of-time framing.)

\subsection{The cosmological acceleration: survives but underdetermined}

The MOND-like scale was predicted as $a_0=cH_0/2\pi$ from a de~Sitter
Gibbons--Hawking $2\pi$, the same factor used for local horizons; the
cosmological-constant alternative $c^2\sqrt\Lambda$ was excluded at $8.4\times$.
But the data do not select the coefficient: SSV's own halo estimator favours
$1/2\pi$ while the independent radial-acceleration relation mildly favours
Verlinde's $1/6$, and the two straddle by $\sim$6\%. The prediction is consistent
and falsifiable (the decisive separator is high-redshift evolution) but
unconfirmed, and the absolute scale is conceded in parallel with $G$.

\section{The diagnosis: form yes, magnitude no}

One pattern runs through every sector. SSV reproduces the \emph{form} of the right
answer cheaply, and concedes the \emph{magnitude}. Time dilation: form yes, $G$
no. Area law: form yes, coefficient no. Hawking temperature: form yes, factor of
four no. Flat rotation curves: form yes, halo amplitude not derived.
Cosmological $a_0$: form yes, scale conceded. Masses: scales handed in, the rest
fitted.

The reason is not bad luck; it is structural, and it is the real lesson. Kinematic
forms --- metric structure, area laws, dispersion relations, conservation laws ---
follow from \emph{locality and symmetry}, which an emergent medium inherits almost
automatically. The dynamical magnitudes --- the coupling constants --- live in the
\emph{microscopic} theory below the effective cutoff. In a superfluid vacuum the
induced gravitational coupling is trans-Planckian by construction (Volovik), so
the effective theory that produces the beautiful forms is structurally the wrong
place to compute the constants. An emergent-gravity model gets geometry for free
and the gravitational constant never. Recognising this early would have saved
years; it is the single most useful thing the programme learned.

\section{The surviving ontology and its honest home}

What remains attractive is the ontology: local time as the rate of phase update of
the medium; a finite update capacity that loading consumes, slowing the local
clock; and gravity as the response of a holographic record to that loading. Each
component already has a mainstream owner --- time-as-phase-rate is the Josephson
relation plus analogue $g_{00}$; finite update capacity is the
Margolus--Levitin/Lloyd bound on the rate of state change; time-as-change is the
relational tradition (Barbour; Connes--Rovelli); screen-update-as-gravity is
Jacobson/Verlinde/Padmanabhan. The ontology adds interpretive coherence (one
phase field carrying the local clock, the static near-field, and the long-range
encoding) but no prediction those programmes lack.

Two honest caveats belong with it. First, an internal tension: the bulk-realist
reading (the medium's local state slows time) and the holographic reading
(``nothing bends space in the bulk; geometry is localised into the encoding'')
are reconciled only by a duality claim that carries no falsifier of its own; it
is a heuristic frame, not a mechanism. Second, the one component that does pull
weight is \emph{presentism}: because the emergent metric is effective and need not
be geodesically complete, there is no real future infinity, which dissolves the
obstacle that breaks global de~Sitter holography and routes naturally to
observer / static-patch holography (Chandrasekaran--Longo--Penington--Witten). But
that removes an obstacle without building the dual: the present cosmological
horizon still has a Gibbons--Hawking entropy whose microstates the medium does not
supply --- the same magnitude gap as everywhere else. The honest home for the
ontology is therefore foundations / philosophy of physics, not a theory of
gravity.

\section{Lessons that generalise}

\begin{enumerate}[leftmargin=1.6em]
\item \textbf{Forms are cheap, constants are dear.} In any emergent/analogue
  programme, expect to reproduce kinematic forms from locality and symmetry, and
  expect the coupling constants to require microscopic input the effective theory
  structurally cannot provide. Treat a reproduced \emph{form} as confirmation of
  consistency, never as a derivation of the \emph{magnitude}.
\item \textbf{Negatives are the durable output.} The scalar-gravity no-go and the
  mode-counting $G$-shortfall will outlast everything positive the programme
  proposed, precisely because a correct negative generalises and a numerical
  coincidence does not.
\item \textbf{Concreteness exposes what abstraction hides.} Abstract entropic
  gravity conceals the magnitude problem; building the screen out of real degrees
  of freedom and counting them makes the $10^{38}$ shortfall undeniable.
  Demonstrating a known problem concretely is itself worth publishing.
\item \textbf{Distinguish derivation from fit, in writing.} A claim-status table
  that forces every number to be labelled \emph{derived}, \emph{coincidence}, or
  \emph{falsified} is the cheapest defence against a coincidence (the mass ladder,
  the $\mathcal{O}(1)$ prefactors) being quietly promoted to a derivation.
\item \textbf{Pre-registration matters most without referees.} A hard
  negative-results rule and pre-registered decision rules are what let an
  independent programme falsify itself rather than accumulate suggestive
  positives. This is the part of SSV most worth copying.
\end{enumerate}

\section{Conclusion}

SSV set out to derive nearly everything from one equation and two inputs. It ends
having derived the \emph{forms} of analogue gravity (which were already known),
having needed three inputs rather than two once $G$ was conceded, and having
produced two clean negatives and one re-usable numerical result. Measured against
its ambition that is a failure; measured against the negative-results rule it
adopted, the failure is the contribution. The lessons --- forms are cheap and
constants are dear, negatives are the durable output, concreteness exposes the
magnitude gap, and self-imposed pre-registration is what makes honest closure
possible --- are not specific to a superfluid vacuum. They are addressed to anyone
who tries to grow gravity from a medium.

\begin{thebibliography}{99}
\bibitem{bbm} I.\ Bialynicki-Birula and J.\ Mycielski, \emph{Nonlinear wave
  mechanics}, Ann.\ Phys.\ \textbf{100}, 62 (1976).
\bibitem{zlosh} K.\,G.\ Zloshchastiev, \emph{Logarithmic nonlinearity in theories
  of quantum gravity: origin of time and observational consequences}, Grav.\
  Cosmol.\ \textbf{16}, 288 (2010).
\bibitem{volovik} G.\,E.\ Volovik, \emph{The Universe in a Helium Droplet}
  (Oxford, 2003).
\bibitem{josephson} B.\,D.\ Josephson, \emph{Possible new effects in
  superconductive tunnelling}, Phys.\ Lett.\ \textbf{1}, 251 (1962).
\bibitem{unruh} W.\,G.\ Unruh, \emph{Experimental black-hole evaporation?},
  Phys.\ Rev.\ Lett.\ \textbf{46}, 1351 (1981).
\bibitem{visser} M.\ Visser, \emph{Acoustic black holes: horizons, ergospheres and
  Hawking radiation}, Class.\ Quantum Grav.\ \textbf{15}, 1767 (1998);
  gr-qc/9712010.
\bibitem{jacobson} T.\ Jacobson, \emph{Thermodynamics of spacetime: the Einstein
  equation of state}, Phys.\ Rev.\ Lett.\ \textbf{75}, 1260 (1995).
\bibitem{jacobson16} T.\ Jacobson, \emph{Entanglement equilibrium and the Einstein
  equation}, Phys.\ Rev.\ Lett.\ \textbf{116}, 201101 (2016).
\bibitem{verlinde} E.\ Verlinde, \emph{On the origin of gravity and the laws of
  Newton}, JHEP \textbf{04}, 029 (2011).
\bibitem{verlinde16} E.\ Verlinde, \emph{Emergent gravity and the dark universe},
  SciPost Phys.\ \textbf{2}, 016 (2017).
\bibitem{padmanabhan} T.\ Padmanabhan, \emph{Thermodynamical aspects of gravity:
  new insights}, Rep.\ Prog.\ Phys.\ \textbf{73}, 046901 (2010).
\bibitem{srednicki} M.\ Srednicki, \emph{Entropy and area}, Phys.\ Rev.\ Lett.\
  \textbf{71}, 666 (1993).
\bibitem{bkls} L.\ Bombelli, R.\,K.\ Koul, J.\ Lee, R.\,D.\ Sorkin, \emph{Quantum
  source of entropy for black holes}, Phys.\ Rev.\ D \textbf{34}, 373 (1986).
\bibitem{peschel} I.\ Peschel, \emph{Calculation of reduced density matrices from
  correlation functions}, J.\ Phys.\ A \textbf{36}, L205 (2003).
\bibitem{casinihuerta} H.\ Casini and M.\ Huerta, \emph{Entanglement entropy in
  free quantum field theory}, J.\ Phys.\ A \textbf{42}, 504007 (2009).
\bibitem{sakharov} A.\,D.\ Sakharov, \emph{Vacuum quantum fluctuations in curved
  space and the theory of gravitation}, Sov.\ Phys.\ Dokl.\ \textbf{12}, 1040
  (1968).
\bibitem{susskinduglum} L.\ Susskind and J.\ Uglum, \emph{Black hole entropy in
  canonical quantum gravity and superstring theory}, Phys.\ Rev.\ D \textbf{50},
  2700 (1994).
\bibitem{margolus} N.\ Margolus and L.\,B.\ Levitin, \emph{The maximum speed of
  dynamical evolution}, Physica D \textbf{120}, 188 (1998).
\bibitem{lloyd} S.\ Lloyd, \emph{Ultimate physical limits to computation}, Nature
  \textbf{406}, 1047 (2000).
\bibitem{barbour} J.\ Barbour, \emph{The End of Time} (Oxford, 1999).
\bibitem{connesrovelli} A.\ Connes and C.\ Rovelli, \emph{Von Neumann algebra
  automorphisms and time--thermodynamics relation in general covariant quantum
  theories}, Class.\ Quantum Grav.\ \textbf{11}, 2899 (1994).
\bibitem{clpw} V.\ Chandrasekaran, R.\ Longo, G.\ Penington, E.\ Witten, \emph{An
  algebra of observables for de~Sitter space}, JHEP \textbf{02}, 082 (2023);
  arXiv:2206.10780.
\end{thebibliography}

\end{document}
